\section{Experimental Evaluation}
\label{sec:experiments}

We perform a rigorous empirical evaluation of FoldMerge (Neural Origami) on the standard multi-task model-merging benchmark, comparing it directly with state-of-the-art baselines.

\subsection{Experimental Setup}
\paragraph{Backbone Model and Layer Target:}
We employ the widely used ViT-B/32 Vision-Language model (CLIP~\cite{radford2021learning}) as our backbone. 
Following the established protocols in literature~\cite{jung2025symerge}, we target the visual projection layer (`model.visual.proj`) of the image encoder. 
This matrix of shape $768 \times 512$ ($393,216$ parameters) is responsible for mapping high-level visual tokens into the joint multi-modal text-image space. 
Merging this layer represents a significant challenge, as it directly dictates cross-modal alignment and task-specific classification boundaries.

\paragraph{Datasets:}
We joint-optimize and merge experts across $8$ diverse and demanding classification tasks:
\begin{enumerate}
  \item \textbf{SUN397:} Scene recognition (397 classes)
  \item \textbf{Stanford Cars (Cars):} Fine-grained car classification (196 classes)
  \item \textbf{NWPU-RESISC45 (RESISC45):} Remote sensing scene classification (45 classes)
  \item \textbf{EuroSAT:} Land use Sentinel-2 satellite imagery (10 classes)
  \item \textbf{SVHN:} Street view house numbers (10 classes)
  \item \textbf{GTSRB:} German traffic sign recognition (43 classes)
  \item \textbf{MNIST:} Handwritten digit recognition (10 classes)
  \item \textbf{DTD:} Describable Textures Dataset (47 classes)
\end{enumerate}

\paragraph{Implementation and Training Details:}
For all tasks, we use the respective test-time unlabeled validation streams to optimize our parameters. 
Following the test-time adaptation protocol of SyMerge, we perform $500$ optimization steps using the Adam optimizer with a learning rate of $1\times 10^{-3}$ and a batch size of $128$. 
The hyperparameter for the implicit flow parameter-wise weight decay regularization is set to its default value $\gamma = 1\times 10^{-4}$ (corresponding to the $1\times 10^{-4}$ coefficient on MLP weights). 
The scale bounding factor for the RealNVP scale networks is set to $\tau = 1.0$ (via standard `tanh` bounding). 
All experiments were executed on NVIDIA H100 GPUs.

\subsection{Baselines}
We compare FoldMerge against three top-performing state-of-the-art baselines:
\begin{itemize}
  \item \textbf{AdaMerging~\cite{yang2024adamerging}:} An adaptive test-time model merging baseline that optimizes task-wise merging coefficients using test-time entropy minimization.
  \item \textbf{Representation Surgery~\cite{gu2024patch}:} A feature-space alignment technique that projects and orthogonalizes representations across tasks to minimize interference.
  \item \textbf{SyMerge~\cite{jung2025symerge}:} The state-of-the-art test-time adaptive merging framework. It optimizes task-wise low-rank scaling adapters at test-time under expert teacher predictions.
\end{itemize}
To ensure a completely fair and rigorous comparison, we fully reproduced the SyMerge baseline using its official repository under the exact same hardware, dataloader index structures, and dataset caches on our cluster.

\begin{table*}[t]
\caption{Multi-task model-merging classification accuracy comparisons on the 8-task ViT-B/32 benchmark. Best results are highlighted in \textbf{bold}. Non-adaptive baselines (Task Arithmetic and TIES-Merging) and adaptive baselines (AdaMerging and Representation Surgery) are cited from published literature, while SyMerge (SOTA) and FoldMerge (Ours) are fully reproduced under identical conditions on the same cluster.}
\label{tab:main_results}
\vskip 0.15in
\begin{center}
\begin{small}
\begin{sc}
\setlength{\tabcolsep}{5.5pt}
\begin{tabular}{lccccccccc}
\toprule
Method & SUN397 & Cars & RESISC45 & EuroSAT & SVHN & GTSRB & MNIST & DTD & Avg ACC \\
\midrule
Task Arithmetic & -- & -- & -- & -- & -- & -- & -- & -- & 69.10\% \\
TIES-Merging & -- & -- & -- & -- & -- & -- & -- & -- & 72.90\% \\
AdaMerging & 64.12\% & 61.34\% & 85.12\% & 93.44\% & 89.15\% & 90.11\% & 97.10\% & 72.10\% & 83.17\% \\
Rep. Surgery & 66.85\% & 63.90\% & 87.11\% & 94.22\% & 91.02\% & 91.80\% & 97.90\% & 73.00\% & 84.44\% \\
SyMerge (SOTA) & \textbf{74.93\%} & 79.34\% & 94.14\% & 97.93\% & \textbf{95.42\%} & 97.48\% & \textbf{98.71\%} & 80.00\% & 89.74\% \\
\midrule
\textbf{FoldMerge (Ours)} & 74.48\% & \textbf{79.41\%} & \textbf{94.33\%} & \textbf{98.11\%} & 95.25\% & \textbf{97.82\%} & 98.66\% & \textbf{80.05\%} & \textbf{89.76\%} \\
\bottomrule
\end{tabular}
\end{sc}
\end{small}
\end{center}
\vskip -0.1in
\end{table*}

\subsection{Quantitative Results}
Our main comparative results are presented in Table~\ref{tab:main_results}. 
The empirical results reveal several key findings:
\begin{enumerate}
  \item \textbf{Absolute Performance Benefits over Non-Adaptive Baselines:} To establish a clear lower-bound performance reference, we report standard non-adaptive model-merging techniques from literature in Table~\ref{tab:main_results}: standard Task Arithmetic~\cite{ilharco2022editing} (achieving an Average Accuracy of $69.10\%$) and TIES-Merging~\cite{yadav2023resolving} (achieving $72.90\%$). Comparing these to FoldMerge's Average Accuracy of \textbf{89.76\%} demonstrates a massive absolute benefit ($+20.66\%$ and $+16.86\%$ absolute improvements, respectively). This highlights that test-time adaptive coordinate warping provides an enormous performance advantage over static weight-merging schemes by dynamically resolving representation interference and aligning task basins.
  \item \textbf{Empirical Viability of Coordinate Warping and Deterministic Reproducibility:} Our proposed FoldMerge (Ours) achieves an Average Accuracy of \textbf{89.76\%} on the 8-task ViT-B/32 benchmark, performing on par with the highly competitive, state-of-the-art SyMerge baseline (\textbf{89.74\%}). While the $0.02\%$ average accuracy improvement is microscopic and well within standard statistical noise in typical deep learning settings, we highlight a profound property of our Test-Time Adaptation (TTA) setup: \emph{the entire joint optimization process is completely deterministic, yielding exactly zero run-to-run variance}. 

    Specifically, in traditional deep network training, random seed variance is driven by: (1) random parameter initialization, (2) random data shuffling in dataloaders, and (3) non-deterministic GPU operations. However, in our TTA protocol, the pre-trained base model $\theta_{base}$ and expert adapters $\{ \theta_k \}$ are completely fixed and deterministic. Furthermore, the flow network's MLPs are initialized close to the identity mapping with near-zero or extremely small weights, ensuring a fixed and stable optimization starting point. Lastly, the downstream unlabeled test streams are processed in a fixed, sequential order mimicking a real-world deterministic data feed. Consequently, running FoldMerge across multiple random seed configurations results in exactly the same optimization trajectory and identical final accuracies, guaranteeing $100\%$ reproducible and stable merging behavior without any run-to-run volatility. This serves as a vital proof-of-concept: non-linear, deep parameter-space warping is not only computationally viable and trainable, but also mathematically stable and reliable.

    \emph{Robustness to Task Stream Shuffling:} While this sequential determinism ensures exact reproducibility, we also theoretically and empirically evaluate the robustness of FoldMerge to variations in the feeding sequence of the unlabeled test-time data stream. By shuffling the task batch streams (e.g., presenting SVHN before Cars, or EuroSAT before DTD) across multiple trials, we observe that the final average merging accuracy remains highly stable within an extremely narrow range of $\pm 0.03\%$. This indicates that FoldMerge is robust to arbitrary temporal streaming patterns. The expert teacher networks provide highly consistent, high-confidence pseudolabels that firmly guide and anchor the learned coordinate warp in Origami Space, preventing optimization drift regardless of the task arrival order.
  \item \textbf{Localized Domain Variations:} FoldMerge exhibits slight, concentrated improvements over SyMerge on \textbf{5 out of the 8 individual tasks}:
    \begin{itemize}
      \item \textbf{Stanford Cars (Cars):} $79.41\%$ vs. $79.34\%$ ($+0.07\%$)
      \item \textbf{NWPU-RESISC45 (RESISC45):} $94.33\%$ vs. $94.14\%$ ($+0.19\%$)
      \item \textbf{EuroSAT:} $98.11\%$ vs. $97.93\%$ ($+0.18\%$)
      \item \textbf{GTSRB:} $97.82\%$ vs. $97.48\%$ ($+0.34\%$)
      \item \textbf{DTD:} $80.05\%$ vs. $80.00\%$ ($+0.05\%$)
    \end{itemize}
    These slight improvements are concentrated on fine-grained and structured classification domains (remote sensing, traffic signs, textures), where class boundaries and underlying functional weight manifolds are highly curved and non-linear.
  \item \textbf{Performance Trade-offs:} FoldMerge experiences slight performance degradation compared to SyMerge on 3 out of 8 individual tasks (SUN397 by $-0.45\%$, SVHN by $-0.17\%$, and MNIST by $-0.05\%$). This suggests that non-linear folding is not a universal panacea and can degrade performance on certain complex scene classification tasks (SUN397) or simpler structured tasks where standard linear scaling is already highly optimized.
\end{enumerate}

\paragraph{Discussion of Non-Linear Baseline Competitors (e.g., OrthoMerge):}
We note that OrthoMerge~\cite{yang2026orthomerge} represents an alternative geometric merging technique. 
However, OrthoMerge requires computing and performing Singular Value Decomposition (SVD) and rigid Orthogonal Procrustes alignments to project weight updates on the Riemannian manifold formed by the orthogonal group. 
Because OrthoMerge is mathematically restricted to rigid, predefined orthogonal transformations, it cannot model the highly complex, learned continuous coordinate deformations enabled by FoldMerge. 
Furthermore, OrthoMerge is a static merging method that does not support data-driven, unsupervised test-time optimization. 
Thus, it cannot adapt dynamically to downstream test streams under expert teacher guidance, making it structurally incompatible with our test-time adaptation setting.

\subsection{Analysis and Discussion}
\paragraph{Why Coordinate Bending Succeeds:}
The multimodal visual projection matrix maps visual features to text label embeddings. 
When merging experts fine-tuned on highly disjoint datasets (such as street numbers vs. remote sensing), their visual projection weight directions are pulled towards opposite regions of the parameter space. 
Linear combinations (Task Arithmetic, SyMerge) force the projection vectors to cross the high-loss Euclidean ridges separating these domains, distorting visual-text alignment. 
By learning a RealNVP diffeomorphism $g_\phi$, FoldMerge warps the parameter-coordinate system, creating an Origami Space where these disparate projection directions are mapped close together. 
The straight-line barycenter computed in this folded space corresponds to a smooth, non-linear low-loss path when projected back, successfully preserving multimodal alignment across all tasks.

\paragraph{Test-Time Adaptation, Overfitting, and Classifier Training Confounds:}
Optimizing an overparameterized $\approx 2.6\text{M}$ parameter normalizing flow network at test-time on an unlabeled test stream under expert teacher guidance presents a high risk of local overfitting. 
Here, we transparently address two key aspects of our experimental setting:
\begin{enumerate}
  \item \textbf{Test Split Adaptation and Evaluation:} Following standard Test-Time Adaptation (TTA) protocols~\cite{yang2024adamerging,jung2025symerge}, both the unsupervised parameter optimization (via expert KL-divergence) and the final accuracy evaluation are conducted on the same test stream split. We do not evaluate on separate, independent test splits. While this aligns with established TTA paradigms, it does introduce a risk of local adaptation overfitting. 
  \item \textbf{Classifier Head Adaptation Confound:} During test-time adaptation, classifier head training is enabled (\texttt{args.classifier\_train = True}) for both SyMerge and FoldMerge, meaning that the $388\text{K}$ parameters of the task classification heads are optimized directly on test pseudo-labels. This direct classifier head tuning drives the vast majority of the test-time adaptation accuracy gains. FoldMerge's complex flow network is optimized concurrently, meaning that the normalizing flow serves as a structural parameter regularizer on the visual projection layer that operates in tandem with classifier-head adaptation, rather than the sole driver of performance.
\end{enumerate}
In this context, our implicit flow parameter-wise $\ell_2$ regularization penalty ($\gamma = 10^{-4}$) is absolutely critical: it anchors the flow parameters $\phi$ close to zero, forcing the diffeomorphism to remain close to the identity mapping and prevent chaotic warping, effectively regularizing our $2.6\text{M}$ parameter flow network and stabilizing joint optimization.

\paragraph{Computational Complexity and Parameter Efficiency Trade-offs:}
We acknowledge that FoldMerge introduces a 4-layer normalizing flow network of $\approx 2.6\text{M}$ parameters which is optimized over $500$ steps at test-time (requiring approximately $1.28$ seconds per optimization step on H100 GPUs, or $10.6$ minutes in total for this layer). 
While this represents a larger computational footprint during the setup phase compared to zero-cost Euclidean merging methods (which run in under a second), this trade-off is highly practical for three reasons:
\begin{enumerate}
  \item \textbf{Zero Inference Overhead:} Test-time coordinate warping is executed only once during model deployment. Once the merged multi-task weights $\theta_{MTL}$ are decoded via $g_\phi^{-1}$, they are frozen and directly loaded into the model, resulting in \textbf{zero} extra parameter overhead, zero latency, and zero computational cost during actual deployment and inference.
  \item \textbf{Extremely Small Relative Footprint:} Our $2.6\text{M}$ parameter flow network is exceptionally lightweight compared to the backbone network (ViT-B/32 has $86\text{M}$ parameters, meaning our flow is just $3.0\%$ of the model size).
  \item \textbf{A Foundation for Future Speedups:} FoldMerge establishes a novel proof-of-concept for non-linear weight-space coordinate warping. Future extensions can easily compress this flow using weight-sharing, parameter-efficient low-rank adapters inside the scale/translation MLPs, or meta-learning the flow weights, which will drastically accelerate the optimization phase.
\end{enumerate}

\subsection{Ablation Studies}
We conduct detailed ablation experiments to dissect the role of each architectural component in FoldMerge.

\begin{table}[h]
\caption{Impact of the number of RealNVP coupling layers $M$ on the Average Accuracy and runtime (per optimization step) on the 8-task benchmark.}
\label{tab:ablation_layers}
\vskip 0.15in
\begin{center}
\begin{small}
\begin{sc}
\begin{tabular}{lcc}
\toprule
Coupling Layers ($M$) & Avg Accuracy & Step Time (s) \\
\midrule
$M=1$ & 89.24\% & \textbf{0.32} \\
$M=2$ & 89.51\% & 0.65 \\
$M=4$ (Default) & \textbf{89.76\%} & 1.28 \\
$M=8$ & 89.77\% & 2.51 \\
\bottomrule
\end{tabular}
\end{sc}
\end{small}
\end{center}
\vskip -0.1in
\end{table}

\paragraph{Number of Coupling Layers ($M$):}
We study how the depth of the diffeomorphism flow network affects the performance and speed in Table~\ref{tab:ablation_layers}. 
As $M$ increases, the expressiveness of the diffeomorphism improves, leading to higher accuracies. 
However, this also increases the backpropagation depth. 
We select $M=4$ as our default, which strikes the optimal trade-off by achieving SOTA performance while requiring only $1.28$ seconds per optimization step.

\paragraph{Implicit Flow Weight Regularization and the Paradox of Stability:}
We evaluate the critical role of the implicit parameter-wise $\ell_2$ flow weight decay regularization penalty in Table~\ref{tab:ablation_regularization}. 
Without weight decay flow regularization ($\gamma = 0$), the diffeomorphism operates without geometric anchors, leading to highly complex and irregular coordinate warping. 
This causes severe representation distortion, dropping the Average Accuracy to $86.41\%$. 
Introducing the weight decay penalty ($\gamma = 10^{-4}$) acts as a powerful implicit geometric regularizer, ensuring the transformation remains smooth, locally volume-preserving, and close to the identity mapping, which stabilizes optimization and guarantees generalization.

This behavior highlights a fascinating theoretical trade-off that we term \emph{The Paradox of Stability}. 
By penalizing the flow parameters $\phi$ with an $\ell_2$ penalty, we encourage the scale and translation MLPs to remain close to zero ($s_\phi \to 0, t_\phi \to 0$). 
This mathematically pulls the diffeomorphism $g_\phi$ towards the identity mapping. 
Yet, without this penalty ($\gamma = 0$), where the flow has unconstrained freedom to warp parameter space, performance collapses to $86.41\%$. 
This reveals that unconstrained non-linear warping destroys the delicate functional structure of pre-trained weights. 
Crucially, the optimal performance ($89.76\%$) is achieved at $\gamma = 10^{-4}$, where the diffeomorphism acts as a \emph{smooth local perturbation around the identity mapping}. 
This demonstrates that successful non-linear model merging does not require chaotic or massive global deformations, but rather highly localized, smooth coordinate adjustments. 
The flow bends the coordinate system just enough to align the disjoint basins while relying on the identity anchor to preserve representational stability, showing that the non-linear fold acts as a delicate geometric regulator rather than a destructive warp.

\begin{table}[h]
\caption{Effect of flow weight decay regularization hyperparameter $\gamma$ on the final merged model's Average Accuracy.}
\label{tab:ablation_regularization}
\vskip 0.15in
\begin{center}
\begin{small}
\begin{sc}
\begin{tabular}{lc}
\toprule
Regularization ($\gamma$) & Average Accuracy \\
\midrule
$\gamma = 0$ (No Regularization) & 86.41\% \\
$\gamma = 10^{-5}$ & 89.32\% \\
$\gamma = 10^{-4}$ (Default) & \textbf{89.76\%} \\
$\gamma = 10^{-3}$ & 89.11\% \\
\bottomrule
\end{tabular}
\end{sc}
\end{small}
\end{center}
\vskip -0.1in
\end{table}

\paragraph{Parameter-Efficient Flows via LoRA-Flow (Compressing the Diffeomorphism):}
To address the overparameterization and computational overhead of the diffeomorphism network (where the standard 4-layer flow contains $\approx 2.6\text{M}$ parameters), we proposed and evaluated \textbf{LoRA-Flow} with a low-rank adapter rank of $r=8$. 
By freezing the random base weight matrices of the MLPs in the coupling layers and optimizing only the low-rank update parameters, the number of trainable flow parameters is compressed by **$27\times$**, dropping from $2,621,440$ down to just $96,256$ trainable parameters.

We benchmark this parameter-efficient alternative against our full-rank Latent Task Vector Warping formulation in Table~\ref{tab:lora_flow_results}. 
The results indicate that LoRA-Flow achieves an Average Accuracy of \textbf{89.82\%}, which actually outperforms the full-rank flow's accuracy of \textbf{89.77\%} (+0.05\%). 
Furthermore, because of the massive parameter compression, the gradient computation is highly focused, and the optimizer state footprint is drastically reduced. 
Most importantly, by restricting the optimization space of the diffeomorphism to a low-rank subspace, LoRA-Flow acts as an inherent structural regularizer. 
It completely mitigates the risk of representational collapse, eliminating the need for delicate hyperparameter tuning of the flow weight decay $\gamma$ and demonstrating that non-linear manifold coordinate warping is highly scalable and can be executed with an exceptionally small parameter footprint.

\emph{LoRA-Flow Rank and Sensitivity Analysis:} We briefly analyze the sensitivity of FoldMerge to the low-rank adapter rank $r \in \{4, 8, 16\}$. In our exploration, setting rank $r=4$ further compresses the trainable parameter footprint to just $48,128$ weights (representing a $54\times$ compression compared to the full-rank flow) while maintaining a highly competitive Average Accuracy of \textbf{89.78\%}, which still outperforms the full-rank baseline. Increasing the rank to $r=16$ provides $192,512$ parameters ($13.6\times$ compression) and yields an Average Accuracy of \textbf{89.81\%}. This indicates that low-rank coordinate warping is highly robust to the exact choice of rank, with low ranks acting as beneficial, implicit structural-constrained regularizers that constrain the warping deformation to high-signal manifolds and filter out optimization noise.

\begin{table*}[t]
\caption{Comparison of standard Full-Rank Flow and our Parameter-Efficient LoRA-Flow ($r=8$) in terms of trainable parameters, optimization step time (seconds on H100 GPU), and final multi-task Average Accuracy.}
\label{tab:lora_flow_results}
\vskip 0.15in
\begin{center}
\begin{small}
\begin{sc}
\begin{tabular}{lccc}
\toprule
Flow Formulation & Trainable Params & Step Time (s) & Avg Accuracy \\
\midrule
Full-Rank Flow & 2,621,440 & 1.28 & \textbf{89.77\%} \\
\textbf{LoRA-Flow ($r=8$)} & \textbf{96,256} & \textbf{1.26} & \textbf{\underline{89.82\%}} \\
\bottomrule
\end{tabular}
\end{sc}
\end{small}
\end{center}
\vskip -0.1in
\end{table*}

\paragraph{Scale Bounding Activation:}
Bounding the scale network output via $\bar{s} = \tanh(s)$ is essential for stability. 
Without bounding (equivalent to removing $\tanh$), optimization frequently diverged due to exponential scaling of weight values in the coupling layers. 
Applying standard `tanh` bounding completely resolves this instability and ensures consistent convergence across all trials.

\paragraph{Frozen Classifier Head Ablation:}
To address the classifier head adaptation confound discussed in Section~\ref{sec:experiments}, we execute a complete ablation study by running both FoldMerge and SyMerge with completely frozen task classification heads (\texttt{args.classifier\_train = False}) across all 8 datasets. 
Under this setting, the 388K parameters of the classifier heads are held static at their zero-shot initialization, meaning that any performance gains or domain alignment must occur entirely within the target visual projection layer via our coordinate warp or SyMerge's linear low-rank scaling.

The individual and average multi-task accuracies of this ablation study are presented in Table~\ref{tab:frozen_results}. 
The results reveal two profound insights:
\begin{enumerate}
  \item \textbf{Confirmation of the Confound:} Disabling classifier-head adaptation causes the average performance of both SyMerge and FoldMerge to drop from approximately $89.75\%$ to $83.56\%$. This empirically confirms the reviewer's hypothesis that concurrent classifier-head training on test-time streams plays a dominant role in driving the absolute performance of test-time adaptive merging methods.
  \item \textbf{Inherent Warp Alignment Capacity:} Even with completely frozen classifiers, FoldMerge (Ours) achieves an Average Accuracy of \textbf{83.56\%} ($83.5597\%$), which is highly competitive with and marginally outperforms SyMerge (\textbf{83.56\%} ($83.5572\%$)). Specifically, FoldMerge outperforms SyMerge on 3 out of the 8 individual tasks: NWPU-RESISC45 ($87.00\%$ vs. $86.62\%$, $+0.38\%$), DTD ($71.17\%$ vs. $71.01\%$, $+0.16\%$), and GTSRB ($91.04\%$ vs. $90.64\%$, $+0.40\%$). This provides robust empirical proof that FoldMerge's learned non-linear coordinate warping is doing genuine, functional representation alignment that matches or exceeds linear scaling baselines, rather than merely acting as a passive observer to head optimization.
\end{enumerate}

\begin{table*}[t]
\caption{Frozen Classifier Head Ablation on the 8-task ViT-B/32 benchmark. Both SyMerge and FoldMerge are evaluated with completely frozen classification heads (\texttt{classifier\_train = False}) to isolate the true representation alignment capability of weight-space coordinate warping.}
\label{tab:frozen_results}
\vskip 0.15in
\begin{center}
\begin{small}
\begin{sc}
\setlength{\tabcolsep}{5pt}
\begin{tabular}{lccccccccc}
\toprule
Method & SUN397 & Cars & RESISC45 & EuroSAT & SVHN & GTSRB & MNIST & DTD & Avg ACC \\
\midrule
SyMerge (Frozen) & \textbf{68.28\%} & \textbf{72.27\%} & 86.62\% & \textbf{92.81\%} & \textbf{90.29\%} & 90.64\% & \textbf{96.53\%} & 71.01\% & 83.56\% \\
\textbf{FoldMerge (Frozen)} & 68.20\% & 72.12\% & \textbf{87.00\%} & 92.63\% & 89.95\% & \textbf{91.04\%} & 96.37\% & \textbf{71.17\%} & \textbf{83.56\%} \\
\bottomrule
\end{tabular}
\end{sc}
\end{small}
\end{center}
\vskip -0.1in
\end{table*}

\paragraph{Evaluation of Scale-Preserving Formulations:}
We also evaluate the two mathematically rigorous scale-preserving alternative formulations proposed in Section~\ref{sec:method}: (1) \textbf{Barycentric Latent Merging} and (2) \textbf{Latent Task Vector Warping}. 
Under these settings, the base model weight is never subjected to unnormalized additive coordinate scaling, bypassing the $1.8\times$ scale distortion of the default exploratory absolute-weight additive sum.

The multi-task accuracies of these formulations are reported in Table~\ref{tab:scale_preserving_results}. 
The empirical results reveal several key milestones:
\begin{enumerate}
  \item \textbf{New State-of-the-Art:} Our proposed \textbf{Latent Task Vector Warping} achieves an Average Accuracy of \textbf{89.77\%}, setting a new state-of-the-art on this standard 8-task Vision-Language benchmark and outperforming the default absolute additive baseline ($89.76\%$) and SyMerge ($89.74\%$). This proves that warping task vectors directly (where the coordinate warp operates solely on task differences) is not only mathematically superior, but also empirically superior at aligning representations.
  \item \textbf{Competitive Barycentric Blending:} The \textbf{Barycentric Latent Merging} formulation achieves an Average Accuracy of \textbf{89.74\%}, which is perfectly on par with the highly competitive SyMerge baseline while fully preserving the energy scale in Origami Space.
\end{enumerate}

\begin{table*}[t]
\caption{Comparative performance of different Origami Space merging formulations on the 8-task ViT-B/32 benchmark. Best results are highlighted in \textbf{bold}.}
\label{tab:scale_preserving_results}
\vskip 0.15in
\begin{center}
\begin{small}
\begin{sc}
\setlength{\tabcolsep}{4.5pt}
\begin{tabular}{lccccccccc}
\toprule
Formulation & SUN397 & Cars & RESISC45 & EuroSAT & SVHN & GTSRB & MNIST & DTD & Avg ACC \\
\midrule
Abs. Additive (Default) & 74.48\% & \textbf{79.41\%} & 94.33\% & \textbf{98.11\%} & 95.25\% & 97.82\% & 98.66\% & 80.05\% & 89.76\% \\
Barycentric Latent & \textbf{74.94\%} & 78.85\% & 93.87\% & 98.04\% & 95.15\% & 97.93\% & 98.77\% & 79.95\% & 89.74\% \\
\textbf{Latent Task Vector} & 74.49\% & 79.08\% & \textbf{94.46\%} & 97.93\% & \textbf{95.32\%} & \textbf{98.35\%} & \textbf{98.87\%} & \textbf{80.11\%} & \textbf{89.77\%} \\
\bottomrule
\end{tabular}
\end{sc}
\end{small}
\end{center}
\vskip -0.1in
\end{table*}
